% titel: Blatt 2, A2: Molare Masse von Verbindungen
% Teil II: der durchgerechnete Weg.

\subsection*{Gegeben und gesucht}
Gefordert: molare Massen von fünf Verbindungen. Gesucht $M$ in $\mathrm{g/mol}$.

\subsection*{Rechnung}
Atommassen (gerundet): $\mathrm{P}\ 30.97,$ $\mathrm{Cl}\ 35.45,$
$\mathrm{K}\ 39.10,$ $\mathrm{Mn}\ 54.94,$ $\mathrm{O}\ 16.00,$
$\mathrm{Cr}\ 52.00,$ $\mathrm{S}\ 32.06,$ $\mathrm{Mg}\ 24.31,$
$\mathrm{H}\ 1.008,$ $\mathrm{C}\ 12.01,$ $\mathrm{N}\ 14.01$.

\begin{itemize}
  \item[a)] Phosphorpentachlorid $\mathrm{PCl_5}$:
        $M = 30.97 + 5\cdot 35.45 = 30.97 + 177.25 = 208.22\ \mathrm{g/mol}$.
  \item[b)] Kaliumpermanganat $\mathrm{KMnO_4}$:
        $M = 39.10 + 54.94 + 4\cdot 16.00 = 39.10+54.94+64.00 = 158.04\ \mathrm{g/mol}$.
  \item[c)] Chrom(III)sulfat $\mathrm{Cr_2(SO_4)_3}$:
        $M = 2\cdot 52.00 + 3\cdot(32.06+4\cdot 16.00)
        = 104.00 + 3\cdot 96.06 = 104.00 + 288.18 = 392.18\ \mathrm{g/mol}$.
  \item[d)] Magnesiumsulfat-heptahydrat $\mathrm{MgSO_4 \cdot 7\,H_2O}$:
        \begin{align*}
          M &= 24.31 + 32.06 + 4\cdot 16.00
              + 7\cdot(2\cdot 1.008 + 16.00) \\
            &= 24.31 + 32.06 + 64.00 + 7\cdot 18.016
            = 120.37 + 126.11 = 246.48\ \mathrm{g/mol}.
        \end{align*}
  \item[e)] Diethylamin $\mathrm{C_4H_{11}N}$:
        $M = 4\cdot 12.01 + 11\cdot 1.008 + 14.01
        = 48.04 + 11.09 + 14.01 = 73.14\ \mathrm{g/mol}$.
\end{itemize}

\begin{ergebnis}
\begin{align*}
  \mathrm{PCl_5} &= 208.22,\quad
  \mathrm{KMnO_4} = 158.04,\quad
  \mathrm{Cr_2(SO_4)_3} = 392.18, \\
  \mathrm{MgSO_4 \cdot 7\,H_2O} &= 246.48,\quad
  \mathrm{C_4H_{11}N} = 73.14
  \quad(\text{alle in }\mathrm{g/mol}).
\end{align*}
\end{ergebnis}

\begin{kontrolle}
Die Summe der Elementbeitraege umkehren: bei $\mathrm{Cr_2(SO_4)_3}$ ergeben
$2\cdot 52 + 3\cdot 32 + 12\cdot 16 = 392$, was mit $392.18$ übereinstimmt. Die
Werte liegen in der erwarteten Größenordnung (Molekülmassen zwischen ca. 70
und 400 g/mol); das Hydrat ist deutlich schwerer als wasserfreies
$\mathrm{MgSO_4}$ ($\approx 120\ \mathrm{g/mol}$), weil $7\,\mathrm{H_2O}$
($126\ \mathrm{g/mol}$) hinzukommen.
\end{kontrolle}

\begin{merksatz}
Molare Masse = Summe aus (Atomanzahl $\times$ Atommassen). Struktur sicher
ablesen (Klammern, Hydratwasser), Elemente durchzaehlen, dann aufaddieren.
\end{merksatz}
